\chapter{Support Vector Machines and Naive Bayes Classifiers}
\label{ch:svm}
\begin{objectives}
After studying this chapter you will be able to:
\begin{itemize}
  \item explain the \textbf{maximum-margin idea} of a linear SVM and the role of \textbf{support vectors};
  \item explain the \textbf{soft margin} and the effect of the parameter $C$;
  \item use the \textbf{dual form} to show that an SVM depends only on dot products between points;
  \item explain the \textbf{kernel trick} and compute \textbf{polynomial} and \textbf{RBF} kernels by hand;
  \item describe \textbf{nonlinear SVM classification} and tune $C$ and $\gamma$;
  \item explain \textbf{SVM regression} with the $\varepsilon$-tube and the $\varepsilon$-insensitive loss;
  \item apply \textbf{Gaussian}, \textbf{multinomial} and \textbf{Bernoulli Naive Bayes} classifiers, with Laplace smoothing, end to end;
  \item answer 2-, 5- and 10-mark university questions on Unit~III.
\end{itemize}
\end{objectives}

\begin{pullquote}[--- Vladimir Vapnik, inventor of the SVM]
Nothing is more practical than a good theory.
\end{pullquote}

This chapter completes Unit~III. It has two halves:

\begin{itemize}
  \item \textbf{Support Vector Machines (SVMs):} linear classification, nonlinear classification with kernels, and regression;
  \item \textbf{Naive Bayes classifiers:} the probabilistic family for text, medical and many other problems.
\end{itemize}

Throughout the chapter you will meet some cartoon characters who help us picture the ideas: a mayor who builds roads, a strict and a relaxed teacher, a detective who loves clues, and a tunnel full of data points.

%====================================================================
\section{Linear SVM Classification}

\subsection{The Idea: Build the Widest Road}

Imagine two colonies of houses, the \textbf{Blue Colony} and the \textbf{Orange Colony}. The mayor wants to build a road between them so that the two colonies are kept apart. Many roads would do the job, but a good mayor builds the \textbf{widest road possible}, so that no house is ever too close to the traffic. The houses that end up right at the \emph{edge} of the road are the ones that decide where the road can go.

\begin{figure}[!htbp]
\centering
\resizebox{0.82\linewidth}{!}{%
\begin{tikzpicture}
\begin{scope}
\clip (-1.4,-1.3) rectangle (7.6,6.6);
\draw[line width=80.5pt,gray!35] (-1,5)--(5,-1);
\draw[white,dashed,very thick] (-1,5)--(5,-1);
\draw[mlred,dashed,thick] (-1,3)--(3,-1);
\draw[mlred,dashed,thick] (-1,7)--(7,-1);
\end{scope}
\foreach \x/\y in {0/0,1/0,0.4/1.4}  \pic[scale=1.5] at (\x,\y) {house=mlblue};
\foreach \x/\y in {4/4,4/3,3.3/4.6} \pic[scale=1.5] at (\x,\y) {house=mlorange};
\pic[scale=1.5] at (1,1) {house=mlblue};
\pic[scale=1.5] at (3,3) {house=mlorange};
\draw[mlpurple,very thick] (1,1) circle (0.42);
\draw[mlpurple,very thick] (3,3) circle (0.42);
\draw[{Stealth}-{Stealth},very thick,mlpurple] (1.3,1.3)--(2.7,2.7);
\node[font=\scriptsize\bfseries,mlpurple,fill=white,inner sep=2pt,rounded corners] at (2.5,1.55) {margin = road width};
\node[font=\small\bfseries,mlblue] at (-0.5,-0.8) {Blue Colony};
\node[font=\small\bfseries,mlorange!80!black] at (5.9,4.5) {Orange Colony};
\node[font=\scriptsize,align=center,fill=white,rounded corners,inner sep=2pt] at (-0.6,2.75) {support vector\\(house on the road edge)};
\draw[->,thick,mlpurple] (-0.15,2.4)--(0.7,1.4);
\node[font=\scriptsize,align=center,fill=white,rounded corners,inner sep=2pt] at (5.9,1.6) {support vector\\(house on the road edge)};
\draw[->,thick,mlpurple] (5.2,2.0)--(3.4,2.85);
\node[font=\scriptsize,fill=white,rounded corners,inner sep=2pt] at (4.6,0.2) {centre line = decision boundary};
\pic[scale=1.6] at (6.7,5.6) {smile};
\bubble[text width=3.7cm]{2.3,5.75}{6.3,5.65}{Build the \textbf{widest} road between the colonies!}
\end{tikzpicture}}
\caption{A cartoon of the maximum-margin classifier: the widest road between two colonies.}\label{fig:road}
\end{figure}

\begin{explain}
\textbf{Houses} are data points: blue houses are one class and orange houses are the other. The \textbf{grey road} is the margin and the \textbf{white dashed centre line} is the decision boundary. The \textbf{red dashed lines} are the edges of the road.

\medskip
The two \textbf{purple circled houses} sit exactly on the edge of the road. They are the \textbf{support vectors}. If you moved or removed any \emph{other} house (without letting it cross the edge) the road would stay exactly where it is; but if you moved a circled house, the whole road would shift. The \textbf{purple arrow} shows the road width, the \textbf{margin}, which the SVM tries to make as large as possible. A wide road means the classifier keeps a safe distance from both classes, so new houses (new data) are less likely to be misclassified.
\end{explain}

\subsection{Definitions}

\begin{definition}[: Support Vector Machine (SVM)]
A \textbf{Support Vector Machine} is a supervised learning algorithm that classifies data by finding the \textbf{maximum-margin hyperplane}: the separating hyperplane $\vect w\cdot\vect x+b=0$ that is \emph{farthest} from the closest training points of both classes. With labels $y_i\in\{-1,+1\}$ the \textbf{hard-margin} problem is
\[
  \min_{\vect w,b}\ \tfrac12\|\vect w\|^2\qquad\text{subject to}\qquad y_i(\vect w\cdot\vect x_i+b)\ge1\ \ \text{for all }i,
\]
and a new point is classified by $\hat y=\operatorname{sign}(\vect w\cdot\vect x+b)$.
\end{definition}

\begin{definition}[: Margin and support vectors]
The \textbf{margin} is the width $2/\|\vect w\|$ of the empty strip between the two parallel hyperplanes $\vect w\cdot\vect x+b=\pm1$. The \textbf{support vectors} are the training points that lie \emph{on} these hyperplanes (or, for a soft margin, inside them or on the wrong side); they alone determine $\vect w$ and $b$.
\end{definition}

In Chapter~4 we solved a small SVM by hand. Here is a quick reminder of its answer, which we will reuse: for the spam data the support vectors were $(1,1)$ with $y=-1$ and $(3,3)$ with $y=+1$, giving $\vect w=(0.5,0.5)$, $b=-2$ and margin $2/\|\vect w\|=2.83$.

\subsection{Soft-Margin SVM: Strict and Relaxed Teachers}

Real data usually overlap, so no road can perfectly separate the colonies. A \textbf{soft margin} lets a few houses stand on the road or even on the wrong side, but each such house is fined.

\begin{definition}[: Soft-margin SVM]
Introduce a \textbf{slack variable} $\xi_i\ge0$ for each point that measures how far it violates the margin, and solve
\[
  \min_{\vect w,b,\vect\xi}\ \tfrac12\|\vect w\|^2+C\sum_{i=1}^{n}\xi_i\qquad\text{s.t.}\qquad y_i(\vect w\cdot\vect x_i+b)\ge1-\xi_i,\ \ \xi_i\ge0 .
\]
The hyperparameter $C>0$ controls the trade-off. Equivalently the SVM minimises the \textbf{hinge loss} $\max(0,1-y_if(\vect x_i))$ plus the penalty $\frac1{2C}\|\vect w\|^2$.
\end{definition}

Think of $C$ as the \textbf{strictness of a teacher}.

\begin{itemize}
  \item \textbf{Large $C$ = strict teacher:} ``No student may break the rules!'' The road becomes narrow and twists to avoid every violation: low bias, \emph{high variance}, risk of overfitting.
  \item \textbf{Small $C$ = relaxed teacher:} ``A few mistakes are fine.'' The road stays wide and smooth: higher bias, \emph{low variance}, risk of underfitting.
\end{itemize}

\begin{figure}[!htbp]
\centering
\begin{tikzpicture}
\begin{groupplot}[group style={group size=2 by 1,horizontal sep=1.1cm,group name=G},width=7.2cm,height=7.2cm,
  xmin=-1,xmax=6,ymin=-1.5,ymax=6.5,grid=both,grid style={gray!15},ticks=none,title style={font=\small\bfseries,align=center}]
\nextgroupplot[title={(a) $C=0.01$: relaxed teacher\\wide road, margin $5.0$}]
\addplot[only marks,mark=*,mlblue,mark size=2pt] coordinates {(0.38,2.25)(0.44,0.97)(1.26,1.20)(0.85,1.30)(2.46,1.39)(1.43,0.97)(0.71,0.63)(2.25,2.95)(2.76,1.70)(2.13,3.18)(1.58,-0.08)(3.02,1.87)(0.63,0.89)(1.04,0.19)(0.49,1.33)(1.46,1.09)(1.54,1.16)(2.03,1.36)(2.34,0.56)(0.82,-0.56)};
\addplot[only marks,mark=triangle*,mlorange,mark size=2.6pt] coordinates {(3.87,5.15)(4.39,4.63)(2.76,3.53)(2.63,4.26)(4.05,3.71)(2.99,3.00)(4.76,2.55)(3.37,5.00)(3.20,2.02)(3.63,2.60)(4.36,4.22)(3.91,4.07)(2.31,4.68)(1.59,3.27)(4.44,4.45)(2.99,1.87)(4.22,1.92)(2.93,3.63)(3.33,5.01)(2.12,4.05)};
\addplot[very thick,mlgreen,domain=-1:6]{(1.408-0.253*x)/0.310};
\addplot[dashed,mlred,thick,domain=-1:6]{(2.408-0.253*x)/0.310};
\addplot[dashed,mlred,thick,domain=-1:6]{(0.408-0.253*x)/0.310};
\nextgroupplot[title={(b) $C=100$: strict teacher\\narrow road, margin $0.53$}]
\addplot[only marks,mark=*,mlblue,mark size=2pt] coordinates {(0.38,2.25)(0.44,0.97)(1.26,1.20)(0.85,1.30)(2.46,1.39)(1.43,0.97)(0.71,0.63)(2.25,2.95)(2.76,1.70)(2.13,3.18)(1.58,-0.08)(3.02,1.87)(0.63,0.89)(1.04,0.19)(0.49,1.33)(1.46,1.09)(1.54,1.16)(2.03,1.36)(2.34,0.56)(0.82,-0.56)};
\addplot[only marks,mark=triangle*,mlorange,mark size=2.6pt] coordinates {(3.87,5.15)(4.39,4.63)(2.76,3.53)(2.63,4.26)(4.05,3.71)(2.99,3.00)(4.76,2.55)(3.37,5.00)(3.20,2.02)(3.63,2.60)(4.36,4.22)(3.91,4.07)(2.31,4.68)(1.59,3.27)(4.44,4.45)(2.99,1.87)(4.22,1.92)(2.93,3.63)(3.33,5.01)(2.12,4.05)};
\addplot[very thick,mlgreen,domain=-1:6]{(13.175-3.322*x)/1.761};
\addplot[dashed,mlred,thick,domain=-1:6]{(14.175-3.322*x)/1.761};
\addplot[dashed,mlred,thick,domain=-1:6]{(12.175-3.322*x)/1.761};
\end{groupplot}
\pic[scale=1.1] at ([shift={(-0.75,-0.75)}]G c1r1.north east) {smile};
\pic[scale=1.1] at ([shift={(-0.75,-0.75)}]G c2r1.north east) {angry};
\end{tikzpicture}
\caption{The same 40 points fitted with a very small and a very large $C$ (real SVM fits).}\label{fig:Ceffect}
\end{figure}

\begin{explain}
Both panels show the \emph{same} 40 houses (blue circles and orange triangles; the two colonies overlap a little in the middle). The \textbf{green line} is the decision boundary and the \textbf{red dashed lines} are the road edges.

\medskip
\textbf{(a) Relaxed teacher ($C=0.01$):} the road is very wide (margin $5.0$) and covers many houses, so many of them are support vectors. The model pays little attention to the individual rule-breakers and looks only at the overall pattern.

\textbf{(b) Strict teacher ($C=100$):} the road is narrow (margin $0.53$) and the line is tilted to squeeze between the points, because every violation is expensive. The model follows the training data closely.
\end{explain}

\medskip
\begin{center}
\begin{tabular}{@{}lcccc@{}}
\toprule
\textbf{$C$} & \textbf{Support vectors} & \textbf{Margin} & \textbf{Train accuracy} & \textbf{Test accuracy (400 pts)}\\
\midrule
0.01 (relaxed) & 32 & 5.00 & 0.925 & 0.953\\
1 & 11 & 1.63 & 0.900 & 0.948\\
100 (strict) & 8 & 0.53 & 0.950 & 0.940\\
\bottomrule
\end{tabular}
\end{center}

\textbf{Read the table:} as $C$ grows, the margin shrinks (5.00 $\to$ 0.53), fewer points are support vectors (32 $\to$ 8) and the \emph{training} accuracy rises slightly ($0.925\to0.950$) while the \emph{test} accuracy drifts down ($0.953\to0.940$): the strict teacher fits the training set better but generalises a little worse. (The differences in accuracy are small here, since the two classes are only mildly overlapping and 400 test points is a small sample; the large change in margin and support vectors is the main lesson.)

\subsection{The Dual Form and Support Vectors}

The SVM optimisation problem can be rewritten in an equivalent \emph{dual} form using Lagrange multipliers $\alpha_i\ge0$, one per training point.

\begin{definition}[: Dual form of the SVM]
\[
  \max_{\vect\alpha}\ \sum_{i}\alpha_i-\frac12\sum_{i}\sum_{j}\alpha_i\alpha_j\,y_iy_j\,(\vect x_i\cdot\vect x_j)
  \quad\text{s.t.}\quad \sum_i\alpha_iy_i=0,\ \ 0\le\alpha_i\le C .
\]
The weights are $\vect w=\sum_i\alpha_iy_i\vect x_i$ and a new point is classified by
\[
  f(\vect x)=\sum_{i}\alpha_i\,y_i\,(\vect x_i\cdot\vect x)+b .
\]
Only the \textbf{support vectors} have $\alpha_i>0$; all other points have $\alpha_i=0$.
\end{definition}

\textbf{Two important facts:}
\begin{enumerate}
  \item The prediction needs only the \emph{support vectors}, which makes it compact.
  \item The data enter \emph{only through dot products} $\vect x_i\cdot\vect x_j$. This is the key to the \textbf{kernel trick} in the next section.
\end{enumerate}

\begin{example}[: The dual form on the spam data]
Support vectors: $(1,1)$ with $y=-1$ and $(3,3)$ with $y=+1$; all other points have $\alpha=0$.

\medskip
\textbf{Solving for $\alpha$:} by symmetry $\alpha_1=\alpha_2=\alpha$. Then
\[
  \vect w=\alpha(-1)(1,1)+\alpha(+1)(3,3)=\alpha(2,2)=(0.5,0.5)\ \Rightarrow\ \alpha=0.25 .
\]
Check $\sum\alpha_iy_i=0.25(-1)+0.25(+1)=0$ \checkmark. The bias is $b=-2$ as before.

\medskip
\textbf{Predicting the new e-mail $(4,1)$} using only dot products with the support vectors:
\[
  f=\underbrace{0.25(-1)(1\cdot4+1\cdot1)}_{=-1.25}+\underbrace{0.25(+1)(3\cdot4+3\cdot1)}_{=+3.75}-2=\mathbf{0.5}>0\ \Rightarrow\ \text{spam}.
\]
This matches the value $f(4,1)=0.5$ found in Chapter~4 directly from $\vect w$.

\medskip
\textbf{Dual objective value:} $\sum\alpha_i-\frac12\sum_{i,j}\ldots=0.5-0.25=0.25=\frac12\|\vect w\|^2$, equal to the primal value: the primal and dual give the same answer.
\end{example}

\subsection{Practical Points}

\begin{itemize}
  \item \textbf{Scale the features.} The margin is measured with distances, so a feature with a huge range (salary in rupees) would dominate one with a tiny range (age in years). Standardise first.
  \item \textbf{Choose $C$ by cross-validation} on a grid such as $10^{-2},10^{-1},\dots,10^{3}$.
  \item \textbf{More than two classes:} SVMs are binary. For $K$ classes use \emph{one-vs-rest} ($K$ classifiers, pick the highest score) or \emph{one-vs-one} ($K(K-1)/2$ classifiers that vote).
  \item \textbf{No probabilities by default.} The SVM output $f(\vect x)$ is a signed distance, not a probability (scikit-learn can fit one afterwards with \texttt{probability=True}).
\end{itemize}

\procon{Linear SVM classification}
{\item Wide margin gives good generalisation
 \item Effective with many features (text)
 \item Only support vectors matter, so it is compact
 \item Convex problem: a single global optimum}
{\item No probabilities by default
 \item Needs scaled features
 \item Sensitive to the choice of $C$
 \item Slow for very large datasets
 \item Hard margin fails if classes overlap}
{\item Spam and news-article classification
 \item Handwritten digit recognition
 \item Gene-expression and cancer classification
 \item Image classification}

%====================================================================
\section{Nonlinear SVM Classification}

\subsection{When a Straight Line Is Not Enough}

Many real problems cannot be solved by a straight line. For instance, if the ``blue'' houses are in the \emph{middle} of a village and the ``orange'' houses are all around them, no straight road can separate them.

\begin{figure}[!htbp]
\centering
\resizebox{\linewidth}{!}{%
\begin{tikzpicture}
% ---------- left panel : 1-D ----------
\node[font=\small\bfseries] at (2.3,3.5) {(a) Original data: 1 feature $x$};
\draw[-{Stealth},thick] (-1.5,0)--(6.2,0) node[below,font=\small]{$x$};
\foreach \x in {-3,-2,-1,0,1,2,3} \draw (2.3+0.8*\x,0.07)--(2.3+0.8*\x,-0.07) node[below,font=\scriptsize]{\x};
\foreach \x in {-1,0,1} \fill[mlblue] (2.3+0.8*\x,0) circle (0.12);
\foreach \x in {-3,-2.5,2.5,3} \node[mlorange,font=\Large] at (2.3+0.8*\x,0) {$\blacktriangle$};
\pic[scale=1.3] at (2.3,2.2) {sad};
\bubble[text width=3.1cm]{2.3,1.15}{2.3,0.9}{No single cut separates them!}
% arrow between panels
\draw[-{Stealth},ultra thick,mlpurple] (6.5,1.5)--(8.3,1.5) node[midway,above,font=\small,align=center]{lift\\$\varphi(x)=(x,x^2)$};
% ---------- right panel : 2-D ----------
\begin{scope}[xshift=9.8cm]
\node[font=\small\bfseries] at (2.3,3.5) {(b) After the lift: 2 features $(x,x^2)$};
\draw[-{Stealth},thick] (-1.5,0)--(6.2,0) node[below,font=\small]{$x$};
\draw[-{Stealth},thick] (2.3,0)--(2.3,3.3) node[right,font=\small]{$x^2$};
\foreach \x/\y in {-1/1,0/0,1/1} \fill[mlblue] (2.3+0.8*\x,0.3*\y) circle (0.12);
\foreach \x/\y in {-3/9,3/9,-2.5/6.25,2.5/6.25} \node[mlorange,font=\Large] at (2.3+0.8*\x,0.3*\y) {$\blacktriangle$};
\draw[very thick,mlgreen] (-1.4,1.0875)--(6.0,1.0875);
\draw[dashed,mlred] (-1.4,0.3)--(6.0,0.3);
\draw[dashed,mlred] (-1.4,1.875)--(6.0,1.875);
\node[font=\scriptsize,mlgreen!50!black,fill=white,inner sep=1pt] at (5.2,1.0875) {a straight line works!};
\pic[scale=1.3] at (-0.2,2.7) {smile};
\end{scope}
\end{tikzpicture}}
\caption{The lifting trick: data that cannot be separated in one dimension become separable after adding the feature $x^2$.}\label{fig:lift}
\end{figure}

\begin{explain}
\textbf{Panel (a):} the data live on a single number line. The blue circles are in the middle ($x=-1,0,1$) and the orange triangles are on both outsides ($x=\pm2.5,\pm3$). Wherever we put \emph{one} cut-point on the line, some blue and some orange end up on the same side. The sad face is right: no straight cut works.

\medskip
\textbf{Panel (b):} we \emph{lift} every point into two dimensions by adding the new feature $x^2$ (the purple arrow). A point $x=3$ becomes $(3,9)$, and $x=1$ becomes $(1,1)$. Now the orange points have large $x^2$ values and the blue points have small ones, so a \emph{horizontal straight line} (green) separates them perfectly, with red dashed lines marking the margin. Going back to the original line, this horizontal cut corresponds to \emph{two} cut-points, $x\approx-1.9$ and $x\approx+1.9$: a \textbf{nonlinear} boundary in $x$ is a \textbf{linear} boundary in the lifted space.
\end{explain}

\begin{keyidea}
\textbf{A linear model in a higher-dimensional feature space is a nonlinear model in the original space.} The SVM can therefore learn curved boundaries by first mapping the data with a function $\varphi$ and then finding the maximum-margin hyperplane there.
\end{keyidea}

\begin{example}[: Solving the lifted problem end to end]
\textbf{Data:} $x=-1,0,1$ are class $-1$ (blue); $x=\pm2.5,\pm3$ are class $+1$ (orange). Map $\varphi(x)=(x,x^2)$.

\begin{center}
\begin{tabular}{@{}c|ccc|cccc@{}}
\toprule
$x$ & $-1$ & $0$ & $1$ & $-2.5$ & $2.5$ & $-3$ & $3$\\
$\varphi(x)=(x,x^2)$ & $(-1,1)$ & $(0,0)$ & $(1,1)$ & $(-2.5,6.25)$ & $(2.5,6.25)$ & $(-3,9)$ & $(3,9)$\\
Class & $-$ & $-$ & $-$ & $+$ & $+$ & $+$ & $+$\\
\bottomrule
\end{tabular}
\end{center}

\textbf{Step 1: closest opposite points.} In the lifted space the blue points $(\pm1,1)$ have $x^2=1$ and the orange points $(\pm2.5,6.25)$ have $x^2=6.25$. These four are the support vectors. The separating line is horizontal, $\vect w=(0,a)$.

\medskip
\textbf{Step 2: margin conditions.} $a(6.25)+b=+1$ and $a(1)+b=-1$. Subtracting, $5.25\,a=2$ so
\[
  a=0.381,\qquad b=-1-0.381=-1.381,\qquad f(x)=0.381\,x^2-1.381 .
\]

\textbf{Step 3: boundary.} $f(x)=0\iff x^2=3.625\iff x=\pm1.90$. The margin is $2/0.381=5.25$ units in the $x^2$ direction.

\medskip
\textbf{Step 4: check the others.} $f(0)=-1.381\le-1$ \checkmark; $f(\pm3)=0.381(9)-1.381=2.05\ge1$ \checkmark.

\medskip
\textbf{Step 5: predict.} $x=2$: $f=0.381(4)-1.381=0.143>0\Rightarrow$ \textbf{orange}. $x=1.5$: $f=0.381(2.25)-1.381=-0.524<0\Rightarrow$ \textbf{blue}.
\end{example}

\paragraph{The cost of lifting.} Adding more and more features (polynomial features of degree 3, 10, \dots) quickly produces an enormous number of columns: $d$ features and degree $p$ give about $\binom{d+p}{p}$ new features. Computing with them directly becomes very slow. The \textbf{kernel trick} avoids this cost.

\subsection{The Kernel Trick}

Remember that the dual SVM uses the data only through dot products $\vect x_i\cdot\vect x_j$. After the lift we would need $\varphi(\vect x_i)\cdot\varphi(\vect x_j)$. A kernel computes that number \emph{without ever building} $\varphi$.

\begin{definition}[: Kernel function]
A \textbf{kernel} is a function $K(\vect a,\vect b)$ that equals a dot product in some (possibly very high-dimensional) feature space: $K(\vect a,\vect b)=\varphi(\vect a)\cdot\varphi(\vect b)$. The \textbf{kernel trick} replaces every dot product $\vect x_i\cdot\vect x_j$ in the SVM by $K(\vect x_i,\vect x_j)$, so the SVM works in the high-dimensional space \emph{at the cost of computing a simple formula}:
\[
  f(\vect x)=\sum_i\alpha_i\,y_i\,K(\vect x_i,\vect x)+b .
\]
\end{definition}

\begin{center}
\begin{tabular}{@{}p{3.1cm}p{5.5cm}p{4.9cm}@{}}
\toprule
\textbf{Kernel} & \textbf{Formula} & \textbf{Parameters}\\
\midrule
Linear & $K(\vect a,\vect b)=\vect a\cdot\vect b$ & none\\
Polynomial & $K(\vect a,\vect b)=(\gamma\,\vect a\cdot\vect b+r)^{d}$ & degree $d$, $\gamma$, $r$\\
RBF (Gaussian) & $K(\vect a,\vect b)=\exp\!\big(-\gamma\|\vect a-\vect b\|^2\big)$ & $\gamma>0$\\
Sigmoid & $K(\vect a,\vect b)=\tanh(\gamma\,\vect a\cdot\vect b+r)$ & $\gamma$, $r$\\
\bottomrule
\end{tabular}
\end{center}

\begin{example}[: The polynomial kernel really equals a dot product]
Let $\vect a=(1,2)$ and $\vect b=(3,1)$ and take the kernel $K(\vect a,\vect b)=(\vect a\cdot\vect b)^2$.

\medskip
\textbf{Using the kernel (one cheap step):} $\vect a\cdot\vect b=3+2=5$, so $K=5^2=\mathbf{25}$.

\medskip
\textbf{Using the explicit feature map} $\varphi(\vect x)=(x_1^2,\ \sqrt2\,x_1x_2,\ x_2^2)$ (three features):
\[
  \varphi(\vect a)=(1,\ 2.828,\ 4),\qquad \varphi(\vect b)=(9,\ 4.243,\ 1),
\]
\[
  \varphi(\vect a)\cdot\varphi(\vect b)=1(9)+2.828(4.243)+4(1)=9+12+4=\mathbf{25}.
\]
Both ways give $25$, but the kernel never needed the three-dimensional vectors.
\end{example}

\begin{example}[: The RBF kernel as a similarity score]
For $\vect a=(1,2)$ and $\vect b=(3,1)$: $\|\vect a-\vect b\|^2=2^2+1^2=5$.
\begin{center}
\begin{tabular}{@{}lccc@{}}
\toprule
$\gamma$ & $0.1$ & $0.5$ & $2$\\
\midrule
$K=e^{-5\gamma}$ & $e^{-0.5}=0.607$ & $e^{-2.5}=0.082$ & $e^{-10}=0.00005$\\
\bottomrule
\end{tabular}
\end{center}
Also $K(\vect a,\vect a)=e^{0}=1$ (identical points are maximally similar). The RBF kernel is a \textbf{similarity between $0$ and $1$} that falls as the points move apart, and $\gamma$ says \emph{how fast}. The feature space of the RBF kernel has infinitely many dimensions, yet each $K$ is one cheap calculation.
\end{example}

\begin{figure}[!htbp]
\centering
\begin{tikzpicture}
\begin{axis}[width=11cm,height=6cm,xlabel={distance $\|\vect a-\vect b\|$},ylabel={similarity $K=e^{-\gamma\,\|\vect a-\vect b\|^2}$},
  xmin=0,xmax=3,ymin=0,ymax=1.05,samples=100,grid=both,grid style={gray!18},
  legend style={at={(0.98,0.97)},anchor=north east,font=\small}]
\addplot[very thick,mlgreen,domain=0:3]{exp(-0.1*x^2)};
\addplot[very thick,mlblue,domain=0:3]{exp(-1*x^2)};
\addplot[very thick,mlred,domain=0:3]{exp(-10*x^2)};
\legend{$\gamma=0.1$ (wide hill),$\gamma=1$,$\gamma=10$ (narrow spike)}
\end{axis}
\end{tikzpicture}
\caption{The RBF kernel as a function of distance, for three values of $\gamma$.}\label{fig:rbfgamma}
\end{figure}

\begin{explain}
The horizontal axis is the \textbf{distance} between two points and the vertical axis is their \textbf{similarity} (1 = identical, 0 = completely unrelated). Each curve is a different $\gamma$.

\begin{itemize}
  \item \textbf{Green ($\gamma=0.1$):} a wide, gentle hill. Even points that are $2$ units apart are still considered fairly similar. Every training point influences a large neighbourhood, so the model is \emph{smooth} (risk of underfitting).
  \item \textbf{Blue ($\gamma=1$):} a medium hill.
  \item \textbf{Red ($\gamma=10$):} a narrow spike. Only points very close together are similar. Each training point influences just its immediate surroundings, so the boundary becomes wiggly and tightly follows the data (risk of overfitting).
\end{itemize}
\end{explain}

\subsection{How an RBF SVM Decides: Hills of Influence}

Here is a cartoon way to see the RBF-kernel SVM. Each support vector builds a \textbf{hill of influence} around itself: a positive hill for an orange support vector (it ``pulls'' nearby points to the orange class) and a negative hill (a pit) for a blue one. The decision value $f(\vect x)$ at any point is the \emph{sum} of all the hills and pits there.

\begin{figure}[!htbp]
\centering
\begin{tikzpicture}
\begin{axis}[width=11.5cm,height=6.5cm,xlabel={position $x$},ylabel={decision value $f(x)$},
  xmin=-1,xmax=5,ymin=-1.3,ymax=1.3,samples=120,axis lines=middle,
  legend style={at={(0.5,1.04)},anchor=south,legend columns=3,font=\small},xtick={1,2,3},xticklabels={blue SV,boundary,orange SV},ytick=\empty]
\addplot[thick,mlorange,domain=-1:5]{exp(-(x-3)^2)};
\addplot[thick,mlblue,domain=-1:5]{-exp(-(x-1)^2)};
\addplot[very thick,black,domain=-1:5]{exp(-(x-3)^2)-exp(-(x-1)^2)};
\legend{orange SV hill,blue SV pit,sum $=f(x)$}
\fill[mlblue] (axis cs:1,0) circle (3pt);
\fill[mlorange] (axis cs:3,0) circle (3pt);
\fill[mlgreen] (axis cs:2,0) circle (3.5pt);
\node[font=\small,mlgreen!50!black,above] at (axis cs:2,0.12) {$f=0$};
\node[font=\small,mlorange!80!black] at (axis cs:4.3,0.75) {$f>0$: orange};
\node[font=\small,mlblue] at (axis cs:-0.2,-0.75) {$f<0$: blue};
\end{axis}
\end{tikzpicture}
\caption{Hills of influence: two RBF support vectors ($\gamma=1$) and the resulting decision value.}\label{fig:hills}
\end{figure}

\begin{explain}
We work in one dimension for clarity, with one blue support vector at $x=1$ and one orange at $x=3$.

\begin{itemize}
  \item The \textbf{orange curve} is the positive hill $e^{-(x-3)^2}$ centred on the orange support vector.
  \item The \textbf{blue curve} is the negative pit $-e^{-(x-1)^2}$ centred on the blue support vector.
  \item The \textbf{black curve} is their sum, the decision value $f(x)$. It is positive on the right (the orange hill wins), negative on the left (the blue pit wins), and crosses \textbf{zero at $x=2$} (green dot), exactly midway, where both influences are equal.
\end{itemize}
In more dimensions it is the same picture: the decision boundary is where the hills and pits cancel, which can be a curve of any shape.
\end{explain}

\subsection{Choosing $C$ and $\gamma$: Real Results}

We trained RBF SVMs on the \texttt{moons} data (two interleaving half-circles with noise): 180 points to train and 120 to test.

\begin{center}
\begin{tabular}{@{}ccccl@{}}
\toprule
\textbf{$\gamma$} & \textbf{Train accuracy} & \textbf{Test accuracy} & \textbf{Support vectors} & \textbf{Diagnosis}\\
\midrule
0.1 & 0.839 & 0.867 & 79 & too smooth: underfits\\
1 & 0.950 & \textbf{0.992} & 54 & good\\
10 & 0.961 & 0.975 & 115 & starting to overfit\\
100 & \textbf{1.000} & 0.850 & 173 & overfits (memorises)\\
\bottomrule
\end{tabular}
\end{center}

\begin{center}
\begin{tabular}{@{}ccc@{\qquad\qquad}ccc@{}}
\toprule
\textbf{$C$} ($\gamma=1$) & \textbf{Train} & \textbf{Test} & \textbf{Kernel} ($C=1$) & \textbf{Train} & \textbf{Test}\\
\midrule
0.1 & 0.928 & 0.950 & Linear & 0.817 & 0.883\\
1 & 0.950 & \textbf{0.992} & Polynomial (degree 3) & 0.944 & 0.983\\
10 & 0.950 & 0.975 & RBF & 0.950 & \textbf{0.992}\\
1000 & 0.967 & 0.933 & & &\\
\bottomrule
\end{tabular}
\end{center}

\textbf{Interpretation.} (i)~With $\gamma=100$ the training accuracy is a perfect $1.000$ but the test accuracy falls to $0.850$: the classic overfitting signature. (ii)~Too small a $\gamma$ underfits. (iii)~Very large $C$ (1000) also overfits ($0.967$ train, $0.933$ test). (iv)~On this curved data the linear SVM is clearly worse ($0.883$) than the polynomial and RBF kernels. (Test accuracies are based on 120 points, so one point is $0.8\%$.)

\begin{tryit}
Reproduce the experiment with \texttt{scikit-learn}:
\begin{lstlisting}
from sklearn.datasets import make_moons
from sklearn.model_selection import train_test_split
from sklearn.preprocessing import StandardScaler
from sklearn.pipeline import make_pipeline
from sklearn.svm import SVC

X, y = make_moons(n_samples=300, noise=0.25, random_state=0)
Xa, Xb, ya, yb = train_test_split(X, y, test_size=0.4,
                                  random_state=0)

for g in [0.1, 1, 10, 100]:
    model = make_pipeline(StandardScaler(),
                          SVC(kernel="rbf", gamma=g, C=1))
    model.fit(Xa, ya)
    print(g, model.score(Xa, ya), model.score(Xb, yb))
\end{lstlisting}
\emph{Challenge:} replace the loop by \texttt{GridSearchCV} over $C\in\{0.1,1,10,100\}$ and $\gamma\in\{0.1,1,10\}$ and report the best pair.
\end{tryit}

\begin{figure}[!htbp]
\centering
\begin{tikzpicture}
\begin{axis}[width=8.5cm,height=8.5cm,axis equal image,xmin=-1.3,xmax=1.3,ymin=-1.3,ymax=1.3,grid=both,grid style={gray!15},
  xlabel={$x_1$},ylabel={$x_2$},legend style={at={(1.04,0.5)},anchor=west,font=\small}]
\addplot[only marks,mark=*,mlblue,mark size=2.4pt] coordinates {(0.16,0.25)(0.38,-0.02)(0.16,0.37)(0.37,-0.19)(-0.19,0.37)(0.47,-0.05)(-0.26,0.33)(0.33,0.31)(-0.37,0.12)(-0.27,-0.32)(0.39,0.06)(-0.20,-0.27)(-0.25,0.27)(0.08,-0.49)(0.20,-0.32)(0.41,0.14)(-0.10,-0.38)(0.08,-0.46)(0.00,0.30)(-0.27,-0.02)};
\addlegendentry{inner ring}
\addplot[only marks,mark=triangle*,mlorange,mark size=3pt] coordinates {(0.61,-0.77)(0.97,-0.39)(-0.74,0.78)(1.00,0.36)(-0.23,-0.88)(0.83,0.71)(-0.41,0.78)(0.10,1.00)(-0.84,0.19)(0.96,-0.24)(-0.90,-0.57)(0.93,-0.55)(-1.04,-0.12)(0.46,0.97)(-0.96,-0.16)(-0.48,-0.90)(-0.16,0.92)(-0.91,0.44)(0.53,-0.72)(-0.95,-0.24)};
\addlegendentry{outer ring}
\addplot[very thick,mlgreen,domain=0:360,samples=100,variable=\t]({0.676*cos(t)},{0.676*sin(t)});
\addlegendentry{RBF boundary}
\end{axis}
\end{tikzpicture}
\caption{An RBF SVM separates two concentric rings that no straight line can (a real fit, 20 points of each class shown).}\label{fig:circles}
\end{figure}

\begin{explain}
The blue circles form an \textbf{inner ring} and the orange triangles an \textbf{outer ring}. A straight line cannot separate them: a linear SVM got only $58\%$ of these points right. The \textbf{RBF SVM} ($\gamma=2$, $C=10$) classified all of them correctly, and its decision boundary (green) is an almost perfect \textbf{circle of radius $\approx0.68$} lying in the gap between the rings. The curved boundary was never specified by us: the kernel allowed the SVM to find it.
\end{explain}

\begin{keyidea}
\textbf{Choosing a kernel:} start with a \emph{linear} kernel (fast, good for many features and text). If the data look nonlinear, try the \emph{RBF} kernel, which is a good default, and tune $C$ and $\gamma$ by cross-validation on \emph{scaled} features. Use a polynomial kernel when you believe interactions between features of a known degree matter.
\end{keyidea}

\procon{Nonlinear (kernel) SVM}
{\item Learns nonlinear boundaries without building features by hand
 \item Flexible: many kernels
 \item Works well in high dimensions
 \item Strong generalisation with the margin}
{\item Needs careful tuning of $C$, $\gamma$ and the kernel
 \item Training cost grows roughly as $n^2$ to $n^3$: slow for $n\gg10^5$
 \item Hard to interpret
 \item Needs feature scaling}
{\item Face and object recognition
 \item Handwritten digit recognition
 \item Protein and DNA classification
 \item Remote-sensing image classification}

%====================================================================
\section{SVM Regression}

\subsection{The Idea: Points Inside the Tunnel Are Free}

The SVM idea also works for \textbf{regression}. In classification we wanted a road that keeps points \emph{off} the road. In SVM regression we reverse it: we fit a \textbf{tunnel} (tube) around the regression line and try to get \textbf{as many points as possible inside}. Points inside the tunnel cost nothing; only points \emph{outside} are penalised.

\begin{definition}[: $\varepsilon$-insensitive loss]
\[
  L_\varepsilon(y,f(\vect x))=\max\big(0,\ |y-f(\vect x)|-\varepsilon\big).
\]
Errors smaller than $\varepsilon$ are ignored (loss $0$); larger errors cost the amount \emph{by which they exceed $\varepsilon$}.
\end{definition}

\begin{definition}[: Support vector regression (SVR)]
\textbf{SVR} fits a function $f(\vect x)=\vect w\cdot\vect x+b$ (or its kernelised version) that is as \emph{flat} as possible while keeping most training points within a tube of half-width $\varepsilon$:
\[
  \min_{\vect w,b,\vect\xi,\vect\xi^*}\ \tfrac12\|\vect w\|^2+C\sum_i(\xi_i+\xi_i^*)
  \quad\text{s.t.}\quad
  \begin{cases}y_i-f(\vect x_i)\le\varepsilon+\xi_i\\ f(\vect x_i)-y_i\le\varepsilon+\xi_i^*\\ \xi_i,\xi_i^*\ge0.\end{cases}
\]
The slack variables $\xi_i,\xi_i^*$ measure how far a point lies above or below the tube. The \textbf{support vectors} are the points on or outside the tube; points strictly inside have no influence. With a kernel the prediction is $f(\vect x)=\sum_i(\alpha_i-\alpha_i^*)K(\vect x_i,\vect x)+b$.
\end{definition}

\begin{figure}[!htbp]
\centering
\begin{tikzpicture}
\begin{axis}[width=11.5cm,height=7cm,xlabel={$x$},ylabel={$y$},xmin=0,xmax=10,ymin=0,ymax=10,grid=both,grid style={gray!15},ticks=none]
\addplot[dashed,mlred,thick,domain=0:10]{0.8*x+0.5+0.8};
\addplot[dashed,mlred,thick,domain=0:10]{0.8*x+0.5-0.8};
\addplot[very thick,mlgreen,domain=0:10]{0.8*x+0.5};
\fill[mlgreen!10,opacity=0.5] (axis cs:0,0.5-0.8+0) -- (axis cs:10,8.5-0.8+0) -- (axis cs:10,8.5+0.8) -- (axis cs:0,0.5+0.8) -- cycle;
\addplot[only marks,mark=*,mlgreen!70!black,mark size=3pt] coordinates {(1,1.5)(3,2.6)(5,4.2)(6,4.8)(7,5.4)(9,7.9)};
\addplot[only marks,mark=square*,mlred,mark size=3pt] coordinates {(4,4.9)(8,8.0)};
\addplot[only marks,mark=diamond*,mlpurple,mark size=4pt] coordinates {(2,2.9)};
\draw[mlred,thick] (axis cs:4,4.5)--(axis cs:4,4.9);
\draw[mlred,thick] (axis cs:8,7.7)--(axis cs:8,8.0);
\node[font=\scriptsize,mlred,anchor=west] at (axis cs:4.1,4.95) {$\xi>0$};
\node[font=\scriptsize,mlred,anchor=west] at (axis cs:8.1,8.1) {$\xi>0$};
\draw[{Stealth}-{Stealth},thick] (axis cs:9.4,7.1)--(axis cs:9.4,8.7);
\node[font=\scriptsize,anchor=east,fill=white,inner sep=1pt] at (axis cs:9.3,8.0) {$2\varepsilon$};
\node[font=\scriptsize,mlgreen!50!black,anchor=north west] at (axis cs:0.2,9.6) {regression line $f(x)$};
\end{axis}
\pic[scale=1.1] at ([shift={(2.6,-1.35)}]current axis.north west) {smile};
\bubble[text width=3.6cm]{$(current axis.north west)+(6.2,-1.0)$}{$(current axis.north west)+(3.1,-1.3)$}{Points inside the tunnel are \textbf{free}!}
\end{tikzpicture}
\caption{Support vector regression: a tube of half-width $\varepsilon$ around the fitted line.}\label{fig:svr}
\end{figure}

\begin{explain}
The \textbf{green line} is the fitted regression function and the \textbf{red dashed lines} are the tube edges, a distance $\varepsilon$ above and below the line (total width $2\varepsilon$).

\begin{itemize}
  \item The \textbf{green dots} lie \emph{inside} the tunnel. Their errors are smaller than $\varepsilon$, so they cost nothing and \emph{do not affect} the line at all (they are not support vectors).
  \item The \textbf{red squares} lie \emph{outside} the tunnel. They are penalised by the distance $\xi$ to the nearest tube edge (the short red segments). They are support vectors.
  \item The \textbf{purple diamond} sits exactly on the tube edge and is also a support vector.
\end{itemize}
A wider tunnel (large $\varepsilon$) includes more points, gives a simpler model and fewer support vectors; a narrow tunnel follows the data more closely.
\end{explain}

\begin{example}[: The $\varepsilon$-insensitive loss]
Let $\varepsilon=0.5$. Three predictions:

\begin{center}
\begin{tabular}{@{}c|ccc|c@{}}
\toprule
 & \textbf{Point 1} & \textbf{Point 2} & \textbf{Point 3} & \textbf{Total}\\
\midrule
True $y$ & 3.0 & 4.1 & 6.5 & \\
Predicted $f$ & 3.3 & 4.0 & 5.2 & \\
Error $y-f$ & $-0.3$ & $0.1$ & $1.3$ & \\
\midrule
$L_\varepsilon=\max(0,|e|-0.5)$ & $0$ & $0$ & $0.8$ & $\mathbf{0.8}$\\
Squared error $e^2$ & $0.09$ & $0.01$ & $1.69$ & $1.79$\\
\bottomrule
\end{tabular}
\end{center}
Points 1 and 2 lie inside the tube and are \emph{ignored}; only point~3 is charged ($1.3-0.5=0.8$). Ordinary squared error would have charged all three points.
\end{example}

\subsection{Real Results}

We fitted noisy data from $y=\sin x$ (60 points: 42 train, 18 test).

\begin{center}
\begin{tabular}{@{}lccc@{}}
\toprule
\textbf{Model} & \textbf{Train $R^2$} & \textbf{Test $R^2$} & \textbf{Support vectors}\\
\midrule
Linear SVR & 0.539 & 0.278 & 39\\
RBF SVR ($C=10,\gamma=1,\varepsilon=0.1$) & \textbf{0.971} & \textbf{0.924} & 26\\
\bottomrule
\end{tabular}
\end{center}

\begin{center}
\begin{tabular}{@{}cccl@{}}
\toprule
\textbf{$\varepsilon$} (RBF SVR) & \textbf{Support vectors} & \textbf{Test $R^2$} & \textbf{Comment}\\
\midrule
0.05 & 35 & 0.914 & narrow tube, many support vectors\\
0.2 & 17 & \textbf{0.942} & good balance\\
0.5 & 4 & 0.841 & tube too wide: underfits\\
\bottomrule
\end{tabular}
\end{center}

\textbf{Read the tables.} A straight line cannot follow the sine curve ($R^2=0.28$ on test data), while the RBF kernel can ($0.92$). Widening the tube from $0.05$ to $0.5$ reduces the number of support vectors from $35$ to $4$, but when the tube is too wide the fit gets worse.

\begin{tryit}
\begin{lstlisting}
import numpy as np
from sklearn.svm import SVR
from sklearn.model_selection import train_test_split

rng = np.random.default_rng(0)
x = np.sort(rng.uniform(0, 6, 60))
y = np.sin(x) + rng.normal(0, 0.15, 60)
xa, xb, ya, yb = train_test_split(x.reshape(-1, 1), y,
                                  test_size=0.3, random_state=0)

svr = SVR(kernel="rbf", C=10, gamma=1, epsilon=0.1).fit(xa, ya)
print(svr.score(xa, ya), svr.score(xb, yb), len(svr.support_))
\end{lstlisting}
\emph{Challenge:} try \texttt{epsilon=0.3} and \texttt{gamma=10} and compare the train and test $R^2$.
\end{tryit}

\begin{center}
\begin{tabular}{@{}p{3.6cm}p{4.9cm}p{5.2cm}@{}}
\toprule
\textbf{Parameter} & \textbf{Meaning} & \textbf{Effect when increased}\\
\midrule
$\varepsilon$ & Half-width of the tube & Fewer support vectors, simpler (flatter) model\\
$C$ & Penalty for points outside the tube & Fits training data more closely (overfits)\\
$\gamma$ (RBF) & Reach of each support vector & More wiggly curve (overfits)\\
\bottomrule
\end{tabular}
\end{center}

\procon{SVM regression (SVR)}
{\item Ignores small errors, so robust to noise
 \item Sparse: only support vectors matter
 \item Nonlinear fits through kernels
 \item Good generalisation}
{\item Three parameters to tune ($\varepsilon$, $C$, $\gamma$)
 \item Slow for very large datasets
 \item Needs feature scaling
 \item No probabilistic prediction interval}
{\item Stock and sales forecasting
 \item Electricity-load and weather prediction
 \item Estimating material or engineering properties
 \item Travel-time estimation}

%====================================================================
\section{Naive Bayes Classifiers}

In Chapter~3 we introduced Bayes' theorem and the Naive Bayes classifier. Here we study the three main \emph{variants}, one for each type of feature, and solve a full text-classification example. First recall the principle.

\begin{definition}[: Naive Bayes classifier]
A \textbf{Naive Bayes classifier} predicts the class
\[
  \hat y=\argmax_{c}\ P(c)\prod_{j=1}^{d}P(x_j\mid c),
\]
where the ``naive'' assumption is that the features are \emph{conditionally independent given the class}. It is a \textbf{generative} classifier: it models how each class generates the data, $P(\vect x\mid c)$, and then uses Bayes' theorem to get $P(c\mid\vect x)$.
\end{definition}

\subsection{The Detective Who Multiplies Clues}

Think of the classifier as a \textbf{detective}. Each feature is a \emph{clue}. The detective does not worry about how the clues are connected (the naive assumption); each clue is allowed to ``speak'' separately, and the detective \textbf{multiplies the votes} of all clues together with the prior suspicion.

\begin{figure}[!htbp]
\centering
\resizebox{\linewidth}{!}{%
\begin{tikzpicture}[card/.style={draw,rounded corners,thick,fill=mlblue!8,draw=mlblue,align=center,font=\small,inner sep=6pt,anchor=west}]
\pic[scale=1.9] at (0,0) {detective};
\bubble[text width=3.4cm]{0.2,2.6}{0,1.0}{Each clue speaks alone. I just \textbf{multiply}!}
\node[card] (c1) at (3.3,2.1) {Clue 1: word ``win''\\spam $0.200$ \quad ham $0.067$};
\node[card] (c2) at (3.3,0.45) {Clue 2: word ``free''\\spam $0.133$ \quad ham $0.067$};
\node[card] (c3) at (3.3,-1.2) {Clue 3: word ``now''\\spam $0.133$ \quad ham $0.133$};
\node[card,fill=mlorange!12,draw=mlorange] (pr) at (3.3,-2.85) {Prior suspicion\\spam $0.5$ \quad ham $0.5$};
\node[draw,rounded corners,very thick,fill=mlpurple!10,draw=mlpurple,align=center,font=\small,inner sep=8pt,anchor=west] (sc) at (10.4,-0.35) {\textbf{Scoreboard}\\[3pt]spam: $0.5\times0.200\times0.133\times0.133=0.00178$\\[2pt]ham: $0.5\times0.067\times0.067\times0.133=0.00030$};
\node[draw,rounded corners,very thick,fill=mlgreen!18,draw=mlgreen,align=center,font=\small\bfseries,inner sep=8pt,anchor=west] (vd) at (12.6,-3.0) {Verdict: SPAM\\(posterior $=85.7\%$)};
\foreach \c in {c1,c2,c3,pr} \draw[-{Stealth},thick,gray!70!black] (\c.east)--++(0.6,0) |- (sc.west);
\draw[-{Stealth},ultra thick,mlgreen] (sc.south)--(vd.north);
\draw[-{Stealth},thick,gray] (1.2,0.3)--(c2.west);
\end{tikzpicture}}
\caption{Naive Bayes as a detective: the evidence of each word is multiplied with the prior.}\label{fig:detective}
\end{figure}

\begin{explain}
The detective is examining the message ``\textbf{win free now}''. The three blue cards are the clues, each giving how likely the word is under \emph{spam} and under \emph{ham} (probabilities estimated from a tiny training set, worked out in detail below). The orange card is the prior: before reading the message we think spam and ham are equally likely. The \textbf{purple scoreboard} multiplies the prior with each clue. Spam scores $0.00178$ and ham only $0.00030$, so the \textbf{verdict is spam}. The word ``win'' is the strongest clue because it is three times more likely in spam; the word ``now'' is no help at all (equal probability), so it simply cancels out.
\end{explain}

\subsection{Gaussian Naive Bayes (continuous features)}

\begin{definition}[: Gaussian Naive Bayes]
For \textbf{continuous} features, $P(x_j\mid c)$ is modelled by a normal density with a \emph{class-specific} mean $\mu_{jc}$ and variance $\sigma_{jc}^2$ estimated from the training data of class $c$:
\[
  P(x_j\mid c)=\frac{1}{\sqrt{2\pi}\,\sigma_{jc}}\exp\!\Big(-\frac{(x_j-\mu_{jc})^2}{2\sigma_{jc}^2}\Big).
\]
\end{definition}

\textbf{Real-life example: ripe or unripe mango.} A fruit seller measures the \emph{sweetness score} and the \emph{colour} of mangoes. From past data: ripe mangoes have sweetness $\sim\mathcal N(6,1)$ and are yellow with probability $0.9$; unripe mangoes have sweetness $\sim\mathcal N(3,1)$ and are yellow with probability $0.2$. Both classes are equally common.

\begin{figure}[!htbp]
\centering
\begin{tikzpicture}
\begin{axis}[width=11.5cm,height=6.4cm,xlabel={sweetness score},ylabel={density},xmin=0,xmax=9,ymin=0,ymax=0.55,samples=120,axis lines=left,ytick=\empty,
  legend style={at={(0.5,1.04)},anchor=south,legend columns=2,font=\small}]
\addplot[very thick,mlgreen,domain=0:9]{exp(-(x-3)^2/2)/sqrt(2*pi)};
\addplot[very thick,mlorange,domain=0:9]{exp(-(x-6)^2/2)/sqrt(2*pi)};
\legend{{unripe $\mathcal N(3,1)$},{ripe $\mathcal N(6,1)$}}
\draw[dashed,thick] (axis cs:4.5,0)--(axis cs:4.5,0.5);
\node[font=\scriptsize,fill=white,inner sep=1pt] at (axis cs:4.5,0.52) {boundary $4.5$};
\fill[mlred] (axis cs:4,0.242) circle (2.5pt); \fill[mlred] (axis cs:4,0.054) circle (2.5pt);
\draw[mlred,thin] (axis cs:4,0) -- (axis cs:4,0.242);
\node[font=\scriptsize,mlred,anchor=east] at (axis cs:3.9,0.30) {$0.242$};
\node[font=\scriptsize,mlred,anchor=west] at (axis cs:4.07,0.08) {$0.054$};
\end{axis}
\pic[scale=0.9] at ([shift={(10.0,4.9)}]current axis.south west) {smile};
\node[font=\small\bfseries,mlorange!80!black] at ([shift={(10.0,4.0)}]current axis.south west) {ripe};
\pic[scale=0.9] at ([shift={(1.6,4.9)}]current axis.south west) {sad};
\node[font=\small\bfseries,mlgreen!60!black] at ([shift={(1.6,4.0)}]current axis.south west) {unripe};
\end{tikzpicture}
\caption{Class-conditional densities of sweetness for ripe and unripe mangoes (Gaussian Naive Bayes in one dimension).}\label{fig:mango}
\end{figure}

\begin{explain}
Each bell curve shows how sweetness is distributed in one class: unripe mangoes cluster around $3$ (green) and ripe ones around $6$ (orange). The curves cross at $4.5$, halfway between the means: with equal priors, a mango sweeter than $4.5$ is classified ripe and a less sweet one unripe. For a mango of sweetness $4$ (red dots) the unripe curve is higher ($0.242$) than the ripe curve ($0.054$), so the model leans towards \emph{unripe}.
\end{explain}

\begin{example}[: Mango classification, one feature]
A mango has sweetness $x=4$. Equal priors $0.5$ each.
\begin{align*}
  P(4\mid\text{ripe})&=\tfrac{1}{\sqrt{2\pi}}\,e^{-(4-6)^2/2}=0.3989\,e^{-2}=0.0540,\\
  P(4\mid\text{unripe})&=\tfrac{1}{\sqrt{2\pi}}\,e^{-(4-3)^2/2}=0.3989\,e^{-0.5}=0.2420 .
\end{align*}
\[
  P(\text{ripe}\mid x=4)=\frac{0.5(0.0540)}{0.5(0.0540)+0.5(0.2420)}=\frac{0.0540}{0.2960}=\mathbf{0.182}\ \Rightarrow\ \text{unripe}.
\]
\textbf{Effect of the prior.} If ripe mangoes were four times as common (priors $0.8$ ripe, $0.2$ unripe): $P(\text{ripe}\mid x=4)=\frac{0.8(0.0540)}{0.8(0.0540)+0.2(0.2420)}=\frac{0.0432}{0.0916}=0.472$, much closer to $0.5$. A strong prior can move the decision.
\end{example}

\begin{example}[: Mango classification, two features (sweetness and colour)]
A mango has sweetness $4.8$ and is \textbf{yellow}. Equal priors.

\medskip
\begin{center}
\begin{tabular}{@{}lcc@{}}
\toprule
\textbf{Quantity} & \textbf{Ripe} & \textbf{Unripe}\\
\midrule
Prior $P(c)$ & $0.5$ & $0.5$\\
$P(\text{sweetness}=4.8\mid c)$ (Gaussian) & $\varphi(-1.2)=0.1942$ & $\varphi(1.8)=0.0790$\\
$P(\text{yellow}\mid c)$ & $0.9$ & $0.2$\\
\midrule
\textbf{Score} $=P(c)\prod P(x_j\mid c)$ & $0.5(0.1942)(0.9)=0.0874$ & $0.5(0.0790)(0.2)=0.0079$\\
\bottomrule
\end{tabular}
\end{center}
\[
  P(\text{ripe}\mid\text{data})=\frac{0.0874}{0.0874+0.0079}=\mathbf{0.917}\ \Rightarrow\ \text{ripe}.
\]
\textbf{Teamwork of the clues:} sweetness $4.8$ alone is a moderately ripe-looking value ($P=0.1942/(0.1942+0.0790)=0.71$), and the yellow colour makes it much more convincing ($0.917$).
\end{example}

\subsection{Multinomial Naive Bayes (word counts)}

\begin{definition}[: Multinomial Naive Bayes]
For \textbf{count data} such as word counts in a document, a class $c$ is described by a probability $P(w\mid c)$ for each word $w$ in the vocabulary of size $V$. With Laplace smoothing $\alpha$,
\[
  P(w\mid c)=\frac{N_{wc}+\alpha}{N_c+\alpha V},
\]
where $N_{wc}$ is the number of times word $w$ appears in training documents of class $c$ and $N_c$ is the total number of word occurrences in class $c$. A document with words $w_1,\dots,w_m$ is scored by $P(c)\prod_{k=1}^{m}P(w_k\mid c)$.
\end{definition}

\textbf{Real-life example: spam or ham SMS.} A tiny training set of four messages:

\begin{center}
\begin{tabular}{@{}lll@{}}
\toprule
\textbf{Message} & \textbf{Class} & \textbf{Words}\\
\midrule
``win cash now'' & spam & win, cash, now\\
``win free prize'' & spam & win, free, prize\\
``meet me now'' & ham & meet, me, now\\
``lunch with me'' & ham & lunch, with, me\\
\bottomrule
\end{tabular}
\end{center}

\begin{example}[: Multinomial Naive Bayes, end to end]
\textbf{Step 1: priors.} 2 spam and 2 ham messages: $P(\text{spam})=P(\text{ham})=0.5$.

\medskip
\textbf{Step 2: counts.} The vocabulary has $V=9$ words: cash, free, lunch, me, meet, now, prize, win, with. Each class has $N_c=6$ word occurrences.

\medskip
\textbf{Step 3: smoothed probabilities} with $\alpha=1$, so every denominator is $6+9=15$:

\begin{center}
\small
\begin{tabular}{@{}l|ccccccccc@{}}
\toprule
\textbf{Word} & win & cash & now & free & prize & meet & me & lunch & with\\
\midrule
Count in spam & 2 & 1 & 1 & 1 & 1 & 0 & 0 & 0 & 0\\
$P(w\mid\text{spam})$ & $\frac3{15}$ & $\frac2{15}$ & $\frac2{15}$ & $\frac2{15}$ & $\frac2{15}$ & $\frac1{15}$ & $\frac1{15}$ & $\frac1{15}$ & $\frac1{15}$\\
\midrule
Count in ham & 0 & 0 & 1 & 0 & 0 & 1 & 2 & 1 & 1\\
$P(w\mid\text{ham})$ & $\frac1{15}$ & $\frac1{15}$ & $\frac2{15}$ & $\frac1{15}$ & $\frac1{15}$ & $\frac2{15}$ & $\frac3{15}$ & $\frac2{15}$ & $\frac2{15}$\\
\bottomrule
\end{tabular}
\end{center}

\medskip
\textbf{Step 4: classify ``win free now''.}
\begin{align*}
\text{score(spam)}&=0.5\cdot\tfrac3{15}\cdot\tfrac2{15}\cdot\tfrac2{15}=0.5\cdot\tfrac{12}{3375}=0.001778,\\
\text{score(ham)}&=0.5\cdot\tfrac1{15}\cdot\tfrac1{15}\cdot\tfrac2{15}=0.5\cdot\tfrac{2}{3375}=0.000296 .
\end{align*}
\[
P(\text{spam}\mid\text{message})=\frac{12}{12+2}=\frac{6}{7}=\mathbf{0.857}\ \Rightarrow\ \textbf{spam}.
\]

\medskip
\textbf{Step 5: classify ``meet me now''.} spam: $0.5\cdot\frac1{15}\cdot\frac1{15}\cdot\frac2{15}\to2$ units; ham: $0.5\cdot\frac2{15}\cdot\frac3{15}\cdot\frac2{15}\to12$ units. So $P(\text{ham})=\frac{12}{14}=0.857$: \textbf{ham}.

\medskip
\textbf{Why smoothing matters.} Without smoothing, $P(\text{win}\mid\text{ham})=\frac06=0$, and any message containing ``win'' would get a ham score of exactly $0$, whatever its other words. With $\alpha=1$ it is $\frac1{15}$: small but not zero.
\end{example}

\subsection{Bernoulli Naive Bayes (word present or absent)}

\begin{definition}[: Bernoulli Naive Bayes]
For \textbf{binary} features $x_j\in\{0,1\}$ (word present or absent), each feature has a probability $p_{jc}=P(x_j=1\mid c)$ and
\[
  P(x_j\mid c)=p_{jc}^{\,x_j}(1-p_{jc})^{\,1-x_j},\qquad p_{jc}=\frac{n_{jc}+\alpha}{N_c^{\text{docs}}+2\alpha},
\]
where $n_{jc}$ is the number of documents of class $c$ containing feature $j$. Unlike the multinomial model, it \emph{explicitly uses the absence of a word} as evidence.
\end{definition}

\begin{example}[: Bernoulli Naive Bayes on the same messages]
Each class has $2$ documents, so $p_{jc}=\frac{n_{jc}+1}{4}$.

\begin{center}
\small
\begin{tabular}{@{}l|ccccccccc@{}}
\toprule
\textbf{Word} & win & cash & now & free & prize & meet & me & lunch & with\\
\midrule
$P(\text{present}\mid\text{spam})$ & $0.75$ & $0.5$ & $0.5$ & $0.5$ & $0.5$ & $0.25$ & $0.25$ & $0.25$ & $0.25$\\
$P(\text{present}\mid\text{ham})$ & $0.25$ & $0.25$ & $0.5$ & $0.25$ & $0.25$ & $0.5$ & $0.75$ & $0.5$ & $0.5$\\
\bottomrule
\end{tabular}
\end{center}

\medskip
Classify ``win free now'': present words win, free, now; \textbf{absent} words cash, prize, meet, me, lunch, with.
\begin{align*}
\text{spam:}\ & 0.5\times\underbrace{0.75\cdot0.5\cdot0.5}_{\text{present}}\times\underbrace{0.5\cdot0.5\cdot0.75\cdot0.75\cdot0.75\cdot0.75}_{\text{absent}}=0.01483,\\
\text{ham:}\ & 0.5\times\underbrace{0.25\cdot0.25\cdot0.5}_{\text{present}}\times\underbrace{0.75\cdot0.75\cdot0.5\cdot0.25\cdot0.5\cdot0.5}_{\text{absent}}=0.00055 .
\end{align*}
$P(\text{spam})=\dfrac{0.01483}{0.01483+0.00055}=\mathbf{0.964}$: spam, with an even higher confidence than the multinomial model, because the \emph{absence} of ham-typical words (me, meet, lunch, with) is also evidence.
\end{example}

\begin{center}
\begin{tabular}{@{}p{2.8cm}p{3.2cm}p{3.6cm}p{3.9cm}@{}}
\toprule
\textbf{Variant} & \textbf{Feature type} & \textbf{Models $P(x_j\mid c)$ by} & \textbf{Typical use}\\
\midrule
Gaussian & Continuous measurements & Normal density per class & Medical measurements, sensors\\
Multinomial & Counts (word counts, frequencies) & Word probabilities with smoothing & Text classification, spam\\
Bernoulli & Binary (present/absent) & Presence probability per feature & Short texts, binary indicators\\
\bottomrule
\end{tabular}
\end{center}

\begin{tryit}
Confirm the hand calculations with \texttt{scikit-learn}:
\begin{lstlisting}
from sklearn.feature_extraction.text import CountVectorizer
from sklearn.naive_bayes import MultinomialNB, BernoulliNB

docs = ["win cash now", "win free prize", "meet me now",
        "lunch with me"]
labels = [1, 1, 0, 0]                      # 1 = spam, 0 = ham

cv = CountVectorizer(binary=False)
mnb = MultinomialNB(alpha=1).fit(cv.fit_transform(docs), labels)
print(mnb.predict_proba(cv.transform(["win free now"])))
# -> [[0.143 0.857]]   P(ham), P(spam)

cvb = CountVectorizer(binary=True)
bnb = BernoulliNB(alpha=1).fit(cvb.fit_transform(docs), labels)
print(bnb.predict_proba(cvb.transform(["win free now"])))
# -> [[0.036 0.964]]
\end{lstlisting}
The outputs, $0.857$ (multinomial) and $0.964$ (Bernoulli) for ``spam'', agree with the hand calculation.
\end{tryit}

\subsection{Generative vs.\ Discriminative Classifiers}

\begin{center}
\begin{tabular}{@{}p{3.4cm}p{5.2cm}p{5.2cm}@{}}
\toprule
 & \textbf{Naive Bayes (generative)} & \textbf{Logistic regression / SVM (discriminative)}\\
\midrule
What is learned & $P(\vect x\mid c)$ and $P(c)$ & $P(c\mid\vect x)$ or the boundary directly\\
Training & Counting; very fast & Optimisation; slower\\
Data needed & Works with little data & Needs more data, but usually more accurate with plenty\\
Feature dependence & Assumed independent & Can use correlated features\\
Probabilities & Often over-confident & Logistic: better calibrated\\
\bottomrule
\end{tabular}
\end{center}

\begin{warning}
Naive Bayes often classifies well even when its probabilities are \emph{wrong}. If the posterior says $0.999$, do not trust that number literally: the independence assumption multiplies correlated evidence several times and makes the model over-confident. Use it to \emph{rank} or \emph{decide}, not to quote exact probabilities.
\end{warning}

\procon{Naive Bayes classifiers}
{\item Extremely fast to train and predict
 \item Works with little training data and many features
 \item Handles multiple classes naturally
 \item Excellent baseline for text}
{\item Independence assumption is rarely true
 \item Probabilities over-confident
 \item Zero-frequency problem (needs smoothing)
 \item Cannot learn feature interactions}
{\item Spam filtering and sentiment analysis
 \item News and document categorisation
 \item Medical diagnosis (Gaussian)
 \item Real-time prediction}

%====================================================================
\section{Comparing the Methods of Unit~III}

\begin{center}
\small
\begin{tabular}{@{}p{3.2cm}p{3.6cm}p{3.4cm}p{3.6cm}@{}}
\toprule
 & \textbf{Kernel SVM} & \textbf{Naive Bayes} & \textbf{Random Forest / Boosting}\\
\midrule
Training speed & Slow for large $n$ & \textbf{Very fast} & Moderate\\
Prediction speed & Fast (few support vectors) & \textbf{Very fast} & Moderate\\
Nonlinear? & Yes (kernel) & Mild & \textbf{Yes}\\
Needs scaling? & \textbf{Yes} & No & No\\
Probabilities? & Not directly & Yes (over-confident) & Yes\\
Works with little data & Yes & \textbf{Yes} & Less so\\
Interpretability & Low & Moderate & Low--moderate\\
\end{tabular}
\end{center}

%====================================================================
\begin{summary}
\begin{itemize}
  \item \textbf{Linear SVM:} the maximum-margin hyperplane; margin $2/\|\vect w\|$; support vectors lie on the road edge and alone determine the model.
  \item \textbf{Soft margin:} slack $\xi_i$ and parameter $C$; large $C$ = strict (narrow margin, overfit), small $C$ = relaxed (wide margin, underfit).
  \item \textbf{Dual form:} $f(\vect x)=\sum\alpha_iy_i(\vect x_i\cdot\vect x)+b$; only support vectors have $\alpha_i>0$; data appear only through dot products.
  \item \textbf{Nonlinear SVM:} lift the data with $\varphi$; the \textbf{kernel trick} replaces dot products by $K(\vect a,\vect b)=\varphi(\vect a)\cdot\varphi(\vect b)$. Polynomial: $(\gamma\vect a\cdot\vect b+r)^d$; RBF: $e^{-\gamma\|\vect a-\vect b\|^2}$. Large $\gamma$ or $C$ overfits.
  \item \textbf{SVR:} $\varepsilon$-tube; loss $\max(0,|y-f|-\varepsilon)$; points inside the tube are free; support vectors are on or outside.
  \item \textbf{Naive Bayes:} $\hat y=\argmax_cP(c)\prod_jP(x_j\mid c)$; Gaussian (continuous), multinomial (counts), Bernoulli (binary); use Laplace smoothing $\frac{N_{wc}+\alpha}{N_c+\alpha V}$.
  \item Worked results: SVM lifted boundary at $x=\pm1.90$; RBF kernel $e^{-0.5\cdot5}=0.082$; SMS ``win free now'' spam with $0.857$ (multinomial) and $0.964$ (Bernoulli); mango ripe with $0.917$.
\end{itemize}
\end{summary}

%====================================================================
\section*{Exam Corner}
\addcontentsline{toc}{section}{Exam Corner}

\subsection*{A.\ Very Short Answer Questions (2 marks each)}
\begin{enumerate}
  \item Define a support vector machine. \qmarks{2}
  \item What are support vectors? What is the margin? \qmarks{2}
  \item What is the role of the parameter $C$ in a soft-margin SVM? \qmarks{2}
  \item Write the decision function of an SVM in dual form. \qmarks{2}
  \item What is the kernel trick? \qmarks{2}
  \item Write the formulas of the polynomial and RBF kernels. \qmarks{2}
  \item What does a large $\gamma$ do in the RBF kernel? \qmarks{2}
  \item Define the $\varepsilon$-insensitive loss. \qmarks{2}
  \item What is the effect of increasing $\varepsilon$ in SVR? \qmarks{2}
  \item State the independence assumption of Naive Bayes. \qmarks{2}
  \item What is Laplace smoothing? \qmarks{2}
  \item Name the three types of Naive Bayes and their feature types. \qmarks{2}
\end{enumerate}

\begin{modelans}[: Q5, Q7 and Q9]
\textbf{Q5.} The kernel trick is the replacement of every dot product $\vect x_i\cdot\vect x_j$ in the SVM by a kernel $K(\vect x_i,\vect x_j)=\varphi(\vect x_i)\cdot\varphi(\vect x_j)$. The SVM then works in a high-dimensional feature space without ever computing the mapping $\varphi$, which makes nonlinear classification possible at low cost.\par
\textbf{Q7.} A large $\gamma$ makes the similarity fall off very quickly with distance: every training point influences only a tiny neighbourhood, producing a very wiggly decision boundary. It reduces bias but increases variance, so it \emph{overfits}.\par
\textbf{Q9.} A larger $\varepsilon$ gives a wider tube; more points fall inside, fewer are support vectors, and the model becomes simpler (flatter). If $\varepsilon$ is too large the model underfits.
\end{modelans}

\subsection*{B.\ Short Answer Questions (5 marks each)}
\begin{enumerate}
  \item Explain the maximum-margin classifier and support vectors with a diagram. \qmarks{5}
  \item Explain the soft-margin SVM. Discuss the effect of $C$. \qmarks{5}
  \item Explain the dual form of the SVM and why it is important for kernels. \qmarks{5}
  \item Why do we need nonlinear SVMs? Explain the lifting idea with an example. \qmarks{5}
  \item Explain the kernel trick. Describe the linear, polynomial and RBF kernels. \qmarks{5}
  \item Explain the effect of $C$ and $\gamma$ on an RBF SVM. \qmarks{5}
  \item Explain SVM regression with the $\varepsilon$-tube. \qmarks{5}
  \item Explain the Gaussian, multinomial and Bernoulli Naive Bayes classifiers. \qmarks{5}
  \item Explain why Laplace smoothing is needed in Naive Bayes with an example. \qmarks{5}
  \item Compare generative and discriminative classifiers. \qmarks{5}
\end{enumerate}

\begin{modelans}[: B6 (outline for 5 marks)]
\textbf{Parameter $C$} (penalty for margin violations): large $C$ $\to$ few violations, narrow margin, complex boundary, \emph{low bias, high variance} (overfit); small $C$ $\to$ wide margin, simple boundary, \emph{higher bias, low variance} (underfit).\par
\textbf{Parameter $\gamma$} (reach of each support vector in $e^{-\gamma\|a-b\|^2}$): large $\gamma$ $\to$ narrow influence, wiggly boundary, overfit; small $\gamma$ $\to$ wide influence, smooth boundary, underfit.\par
\textbf{Example:} on the moons data ($\gamma=0.1,1,100$) test accuracy was $0.867,\ 0.992,\ 0.850$ while training accuracy rose to $1.0$ at $\gamma=100$. Both parameters are chosen by cross-validation on scaled features.
\end{modelans}

\subsection*{C.\ Long Answer Questions (10 marks each)}
\begin{enumerate}
  \item Explain linear SVM classification: the margin, the optimisation problem, support vectors, soft margin and the dual form. Solve a small numerical example. \qmarks{10}
  \item Explain nonlinear SVM classification. Describe the kernel trick with the polynomial and RBF kernels, including a numerical verification that $K(\vect a,\vect b)=\varphi(\vect a)\cdot\varphi(\vect b)$. \qmarks{10}
  \item Explain SVM regression: the $\varepsilon$-insensitive loss, the optimisation problem, support vectors and the roles of $\varepsilon$, $C$ and $\gamma$. \qmarks{10}
  \item Explain the Naive Bayes classifier. Classify a message using the multinomial model with Laplace smoothing for a given training set. \qmarks{10}
  \item Describe the Gaussian, multinomial and Bernoulli Naive Bayes classifiers with formulas, examples, advantages and disadvantages. \qmarks{10}
  \item Compare SVM, Naive Bayes and ensemble methods (Random Forest / boosting) for classification. \qmarks{10}
\end{enumerate}

\begin{modelans}[: C2 (numerical verification)]
Take $\vect a=(1,2)$, $\vect b=(3,1)$ and $K(\vect a,\vect b)=(\vect a\cdot\vect b)^2$.\par
\emph{Via the kernel:} $\vect a\cdot\vect b=1(3)+2(1)=5$, $K=25$.\par
\emph{Via the feature map} $\varphi(\vect x)=(x_1^2,\sqrt2x_1x_2,x_2^2)$: $\varphi(\vect a)=(1,2\sqrt2,4)$, $\varphi(\vect b)=(9,3\sqrt2,1)$; $\varphi(\vect a)\cdot\varphi(\vect b)=9+2\sqrt2\cdot3\sqrt2+4=9+12+4=25$.\par
Both give $25$, so the kernel equals the dot product in the 3-D feature space while needing only one 2-D dot product.
\end{modelans}

\subsection*{D.\ Numerical Problems}
\begin{enumerate}
  \item An SVM has $\vect w=(2,1)$, $b=-4$. Classify $(1,1)$ and $(3,2)$ and find the margin. \textit{(Answer: $f(1,1)=-1\Rightarrow-1$ (on the margin); $f(3,2)=4\Rightarrow+1$; margin $=2/\sqrt5=0.894$.)}
  \item In 1-D the support vectors are $x=1$ ($y=-1$) and $x=3$ ($y=+1$). Find $w,b$, the margin and the dual values $\alpha_i$. \textit{(Answer: $w+b=-1$, $3w+b=1\Rightarrow w=1$, $b=-2$; margin $=2$; $w=\sum\alpha_iy_ix_i=3\alpha-\alpha=2\alpha=1\Rightarrow\alpha=0.5$ each; check $\sum\alpha_iy_i=0$.)}
  \item Using the dual weights $\alpha=0.5$ at $x=1$ ($y=-1$) and $x=3$ ($y=+1$) with $b=-2$, compute $f(2.5)$. \textit{(Answer: $0.5(-1)(1\cdot2.5)+0.5(1)(3\cdot2.5)-2=-1.25+3.75-2=0.5>0$.)}
  \item Compute the polynomial kernel $K(\vect a,\vect b)=(\vect a\cdot\vect b+1)^2$ for $\vect a=(1,2)$, $\vect b=(2,1)$. \textit{(Answer: $\vect a\cdot\vect b=4$; $K=25$.)}
  \item Compute the RBF kernel for $\vect a=(1,2)$, $\vect b=(2,1)$ with $\gamma=0.5$. \textit{(Answer: $\|\vect a-\vect b\|^2=2$; $K=e^{-1}=0.368$.)}
  \item With $\varphi(x)=(x,x^2)$ compute $\varphi(2)\cdot\varphi(3)$ and the matching kernel $K(2,3)=2\cdot3+(2\cdot3)^2$. \textit{(Answer: $\varphi(2)\cdot\varphi(3)=(2,4)\cdot(3,9)=6+36=42$; $K=6+36=42$.)}
  \item In SVR with $\varepsilon=0.3$: $y=(2.0,3.0,5.0)$, $f=(2.2,3.5,4.1)$. Find the total $\varepsilon$-insensitive loss. \textit{(Answer: errors $-0.2,-0.5,0.9$; losses $0,0.2,0.6$; total $0.8$.)}
  \item Gaussian NB: ripe $\sim\mathcal N(6,1)$, unripe $\sim\mathcal N(3,1)$, equal priors. Find $P(\text{ripe}\mid x=4)$. \textit{(Answer: densities $0.0540$ and $0.2420$; posterior $0.182$.)}
  \item Using the SMS example, classify ``win now'' with multinomial NB. \textit{(Answer: spam $0.5\cdot\frac3{15}\cdot\frac2{15}=0.01333$; ham $0.5\cdot\frac1{15}\cdot\frac2{15}=0.00444$; $P(\text{spam})=0.75$.)}
  \item A word never occurs in the 200 word-tokens of a class; the vocabulary has 1000 words. Find its Laplace-smoothed probability ($\alpha=1$). \textit{(Answer: $\frac{0+1}{200+1000}=0.00083$.)}
\end{enumerate}

\subsection*{E.\ Multiple-Choice Questions}
\begin{enumerate}
  \item The support vectors of an SVM are: \\
  (a) all training points \quad (b) points on or inside the margin \quad (c) the outliers only \quad (d) the test points
  \item The margin of a linear SVM equals: \\
  (a) $\|\vect w\|$ \quad (b) $1/\|\vect w\|$ \quad (c) $2/\|\vect w\|$ \quad (d) $\|\vect w\|^2$
  \item A very large $C$ in a soft-margin SVM tends to: \\
  (a) widen the margin \quad (b) narrow the margin and overfit \quad (c) remove all support vectors \quad (d) make the model linear
  \item In the dual form the data appear only through: \\
  (a) squares \quad (b) dot products \quad (c) means \quad (d) labels
  \item The RBF kernel is: \\
  (a) $\vect a\cdot\vect b$ \quad (b) $(\vect a\cdot\vect b+1)^2$ \quad (c) $e^{-\gamma\|\vect a-\vect b\|^2}$ \quad (d) $\tanh(\vect a\cdot\vect b)$
  \item A very large $\gamma$ in an RBF SVM leads to: \\
  (a) underfitting \quad (b) overfitting \quad (c) a linear boundary \quad (d) fewer features
  \item In SVR, points strictly inside the $\varepsilon$-tube: \\
  (a) are support vectors \quad (b) have no effect on the model \quad (c) have the largest loss \quad (d) are removed from the data
  \item Increasing $\varepsilon$ in SVR typically: \\
  (a) increases support vectors \quad (b) decreases support vectors \quad (c) changes nothing \quad (d) adds features
  \item Naive Bayes assumes features are: \\
  (a) identical \quad (b) conditionally independent given the class \quad (c) all Gaussian \quad (d) uncorrelated with the class
  \item Which variant suits word counts? \\
  (a) Gaussian \quad (b) Multinomial \quad (c) Linear \quad (d) Sigmoid
  \item Laplace smoothing prevents: \\
  (a) overfitting of trees \quad (b) zero probabilities \quad (c) scaling problems \quad (d) ties
  \item Which classifier is \emph{generative}? \\
  (a) SVM \quad (b) logistic regression \quad (c) Naive Bayes \quad (d) k-NN
\end{enumerate}
\textbf{Answers:} 1-b,\ 2-c,\ 3-b,\ 4-b,\ 5-c,\ 6-b,\ 7-b,\ 8-b,\ 9-b,\ 10-b,\ 11-b,\ 12-c.

\subsection*{F.\ Practice Problems}
\begin{enumerate}
  \item Using the 1-D lifted example, classify $x=-2$ and $x=0.5$ with $f(x)=0.381x^2-1.381$. \textit{(Answer: $f(-2)=0.143\Rightarrow$ orange; $f(0.5)=-1.286\Rightarrow$ blue.)}
  \item Draw the margin and support vectors for the data $(1,1),(2,2)$ in class $-1$ and $(4,4),(5,5)$ in class $+1$. Find $\vect w$, $b$ and the margin. \textit{(Hint: closest pair $(2,2)$ and $(4,4)$.)}
  \item \textbf{Coding:} train \texttt{SVC} with linear, polynomial and RBF kernels on \texttt{make\_circles} and report the accuracies.
  \item \textbf{Coding:} use \texttt{GridSearchCV} to tune $C$ and $\gamma$ of an RBF SVM on the \emph{breast cancer} dataset (after scaling).
  \item \textbf{Coding:} write your own spam examples and train \texttt{MultinomialNB} with \texttt{alpha=1} and \texttt{alpha=0.01}; compare predicted probabilities.
\end{enumerate}

\begin{revision}
\begin{center}\small
\begin{tabular}{@{}p{3.6cm}p{9.6cm}@{}}
\toprule
\textbf{Topic} & \textbf{One-line recall}\\
\midrule
Linear SVM & Max-margin hyperplane; margin $2/\|w\|$; min $\frac12\|w\|^2$ s.t.\ $y_i(w\cdot x_i+b)\ge1$.\\
Support vectors & Points on / inside the margin; only they determine $w,b$.\\
Soft margin & Slack $\xi_i$, penalty $C$: large $C$ strict (overfit), small $C$ relaxed (underfit).\\
Dual form & $f(x)=\sum\alpha_iy_i(x_i\cdot x)+b$; $\alpha_i>0$ only for support vectors.\\
Kernel trick & Replace dot product by $K(a,b)=\varphi(a)\cdot\varphi(b)$.\\
Kernels & Linear $a\cdot b$; poly $(\gamma a\cdot b+r)^d$; RBF $e^{-\gamma\|a-b\|^2}$.\\
$C$, $\gamma$ & Large $C$ / large $\gamma$: overfit; small: underfit.\\
SVR & $\varepsilon$-tube; loss $\max(0,|y-f|-\varepsilon)$; points inside are free.\\
Naive Bayes & $\arg\max_cP(c)\prod_jP(x_j\mid c)$; generative; conditional independence.\\
Variants & Gaussian (continuous), multinomial (counts), Bernoulli (binary).\\
Smoothing & $\frac{N_{wc}+\alpha}{N_c+\alpha V}$ avoids zero probabilities.\\
\bottomrule
\end{tabular}
\end{center}
\end{revision}
